\chapter{ATIFLAR, ALINTILAR VE DİPNOTLAR}
Bu bölümde atıflar, alıntılar ve dipnotların nasıl olması gerektiği hakkında bilgi verilecektir.

\section{Atıflar (kaynakların metin içinde gösterimi)}
\subsection{Yazar soyadına göre atıf verme}
Kaynaklar metin içinde yazar soyadı ve tarih belirtilerek verilir. Kaynaklar sayfasında yazar soyadına göre alfabetik olarak sıralama yapılır. Metin içinde kaynak, cümlenin başlangıcında veya içinde verilecekse, Baran (2003) şeklinde, kaynak cümle sonunda verilecekse (Baran, 2003). şeklinde gösterilir. Nokta işareti kaynaktan hemen sonra konulur.

Kaynak birden fazla yazara ait olduğunda, yazar sayısı iki ise, cümle başında veya içinde Uyar ve Yılmaz (2020) şeklinde, cümle sonunda ise (Uyar ve Yılmaz, 2020). şeklinde yazılır. Yazar sayısı ikiden fazla ise cümle başında veya içinde Uyar ve ark. (2018) şeklinde, cümle sonunda ise (Uyar ve ark, 2018). şeklinde yazılır.

Aynı yazara ait ve aynı yıl içinde yayınlanmış yayınlar Uyar (2008a), Uyar (2008b) şeklinde numaralandırılır. Aynı parantez içerisinde aynı yazarın 2 ve daha fazla eserine atıfta bulunma; son yayınlanan eseri en son belirterek aynı parantez içerisinde gösterilebilirler. Örneğin; Son araştırmalar \ldots{}. göstermiştir (Uyar, 2008, 2020).

Aynı parantez içerisinde 2 ya da daha fazla esere atıf, yazar soyadlarına göre alfabetik sıraya uygun şekilde verilir; (Ateş, 2021; Çelik, 2012; Yılmaz, 2022).

İkincil kaynak drumunda (atıf yapılan kaynak başka bir kaynağa atıfta bulunuyorsa), metinde orijinal kaynağa atıfta bulunulur ve parantez içerisinde orijinal kaynağa atıfta bulunan yazara gönderme yapılır. Referans listesinde sadece orijinal kaynağa atıfta bulunan kaynak için giriş yapılır; orijinal kaynak için referans girişi yapılmaz. Kaynaklar bölümünde künyeler yazar soyadına göre sıralanır.

\subsection{Numara ile atıf verme}
Metin içinde [ ] köşeli parantez içinde numaralandırılır. Tezde ilk verilen kaynak [1] numara ile başlar ve veriliş sırasına göre numaralandırılır. Kaynaklara metin içerisinde aşağıdaki biçimlerde atıf yapılır.

[1]\quad{}\quad{}1 nolu kaynak,

[3-7]\quad{}\quad{}3 ve 7 arası (3, 4, 5, 6 ve 7 nolu ) kaynaklar,

[3, 6]\quad{}\quad{}3 ve 6 nolu kaynaklar,

[3, 5, 8]\quad{}3, 5 ve 8 nolu kaynaklar,

[3, 6-9]\quad{}3 ve 6 ile 9 nolu kaynaklar arasındaki kaynaklar.

Aynı isimli birden fazla cildi olan kaynakların, kullanılan her bir cildine ayrı kaynak numarası verilmelidir.

\subsubsection{Kaynakların listelenmesi}
Kaynaklar bölüm başlığı 1 satır aralığı kullanılarak diğer birinci başlıklardaki gibi (önce 72, sonra 18 punto) aralık bırakılarak yazılmalıdır. Kaynaklar ise tek satır aralıklı, önce ve sonra 6 punto (nk) boşluk bırakılarak yazılır.

Tezde atıf yapılan bütün kaynaklar ‘Kaynaklar’ bölümünde listelenmelidir.

Yazar soyadına göre atıf yapılan tezlerde kaynaklar listesi alfabetik olarak sıralanır. Numaralandırma yöntemine göre atıf yapılan tezlerde ise kaynaklar tezde veriliş sırasına göre numaralanır ve sıralanır.

Kaynaklar bölümünde aynı yazarın birden çok eseri yer alıyorsa, eserler, yayın yılına göre eskiden yeniye doğru sıralanır. Aynı yazarın iki farklı eserinin yayın tarihleri aynıysa, kaynaklar listesindeki sıralama, künyede bir sonraki öğe olan eser adına göre yapılır. Atıflarda ayrımı sağlamak için tarih bilgisinin yanına a’dan başlayan harfler eklenir.

Aynı yazarın tek yazarlı eserleri çok yazarlı eserlerinden önce sıralanır. Çok yazarlı iki eserin ilk yazarları aynıysa, ikinci yazarın soyadı, ikinci yazarlar da aynıysa üçüncü yazarın soyadı alfabetik düzende belirleyicidir.

Aynı soyadını taşıyan iki farklı yazarın eserleri adlarına göre alfabetik sıraya girer.

Yazar bir tüzel kuruluşsa, eser, yazar konumundaki tüzel kuruluşun adıyla alfabetik listeye girer.

Bir eserin yazarı veya editörü yoksa künye eser adına hazırlanacağı için, kaynak, eser adından alfabetik listeye girer. Eser adı rakamla başlıyorsa sıralamada rakamın okunuşu dikkate alınır.

İkincil kaynak durumunda (atıf yapılan kaynak başka bir kaynağa atıfta bulunuyorsa), metinde orijinal kaynağa atıfta bulunulur ve parantez içerisinde orijinal kaynağa atıfta bulunan yazara gönderme yapılır. Referans listesinde sadece orijinal kaynağa atıfta bulunan kaynak için giriş yapılır; orijinal kaynak için referans girişi yapılmaz. Kaynaklar bölümünde künyeler yazar soyadına göre sıralanır.

İkincil kaynak, ilgili kaynağa ulaşılamadığı durumlarda “aktaran” ifadesi ile yapılmalıdır. Örneğin;

Parantez içerisinde kullanımı:

(Baran, 1945; akt. Çoruhlu, 2015).

Metin içi kullanımı:

Baran (aktaran Çoruhlu, 2015) 1945’ te \ldots{}

\section{Alıntılar}
Genel olarak alıntılar kelime, imla ve noktalama bakımından aslına uygun olarak yapılır. Alıntı yapılan parçada bir yanlış varsa, doğrusu köşeli parantez içerisinde belirtilmek koşuluyla metin aynen nakledilir.

Kırk kelimeden daha az uzunluktaki kısa alıntılar çift tırnak içerisinde verilir. Alıntının sonunda ilgili kaynağa atıf yapılıp atıftan sonra nokta koyulur.

Kırk kelimeden fazla olan uzun alıntılar tırnak içerisinde gösterilmezler. Uzun alıntılar soldan 1 sekme (1,27 cm) içerden verilir. İçerden verilen uzun alıntılarda, 2 yazı karakteri daha küçük karakter kullanılır. Eğer alıntı yapılan metinde aynı cümle içinde kelime atlanarak yazılıyorsa bu üç nokta (...) ile, eğer bazı cümleler atlanıyorsa bu dört nokta (....) ile belirtilmelidir. Ancak, çok sık ve çok uzun alıntılardan kaçınılması tavsiye edilir. Kısa alıntılardan farklı olarak noktalama atıftan sonra değil de önce yapılır. Örneğin; . (Yılmaz, 2021, s. 196)

40 kelimeden fazla olan alıntı örneği;

Ana metin ana metin ana metin ana metin ana metin ana metin ana metin ana metin ana metin ana metin. Bu görüşe karşı olan bazı çalışmalarda bulunmaktadır:

\begin{alinti}
Bilim, araştırma ve geliştirme faaliyetleri ekonomik kalkınma ve sosyal gelişmenin gerçekleşmesinde etkin araçlar olarak belirmektedir. Bilimsel ve teknolojik araştırmaların üretim ve karar verme mekanizmalarına olan etkileri özellikle gelişmiş ülkelerde hemen her alanda kendini belli etmektedir. Oysa az gelişmiş ülkelerde bu ilişkiyi fark etmek bile zordur. (Eryaman, 1973, 685)
\end{alinti}

Betimleme yöntemi, Kaptan (1995) tarafından şu biçimde tanımlanmıştır:

\begin{alinti}
Olayların, objelerin, varlıkların, kurumların, grupların ve çeşitli alanların ne olduğunu betimlemeye, açıklamaya çalışan incelemelerdir. Bunlar nedir? sorusuna cevap bulmaya yöneliktir. Bununla mevcut durumlar, koşullar, özellikler aynen ortaya konmaya çalışılır. Betimleme araştırmaları, mevcut olayların daha önceki olay ve koşullarla ilişkilerini de dikkate alarak, durumlar arasındaki etkileşimi açıklamayı hedefler (s. 59).
\end{alinti}

Bazıları bu fikre karşı çıktılar:

\begin{alinti}
Aynı anda, aynı yerde bulunmak bütün grup üyeleri arasında samimi bir etkileşim olmasını gerektirmez. Bir ritüeli gerçekleştirmek ya da bir olayı kutlamak için aynı yerde, bir araya toplanan yüzlerce, hatta binlerce kişiyi düşünün. Bu ortamlarda, katılımcılar grubun görsel tezahürünü yaşayıp, fiziksel olarak bir arada olsa da çevrelerindeki kişilerle doğrudan, samimi bir bağlantı kurma şansları grubun büyüklüğüne bağlıdır (Purcell, 1997, ss. 111-112).
\end{alinti}

40 kelimeden uzun olan, kelime ve cümle atlanarak kullanılarak yazılan alıntı örneği:

\begin{alinti}
Bilim, dünyamızda olup biten olguları betimleme ve açıklama yoluyla anlama girişimidir. Olguları betimleme, onları saptama, sınıflama ve dile getirme gibi işlemleri kapsar. .... Bilimin açıklama yönüne gelince bu konu bizi çok daha geniş sorunlara götürecek niteliktedir. ... Hipotez, doğa yasası, teori, nedensellik ve olasılık ilkeleri gibi kavramları ele almaya ihtiyaç vardır (Yıldırım, 1991, 95).
\end{alinti}

Cümle başındaki alıntı örnekleri;

Orhan (2018)’e göre, “Bu tür çalışmalar çok sayıda fiziksel, kimyasal ve biyolojik parametrenin çok sayıda ölçüm noktasında incelenmesini gerektirmektedir” (s. 199).

Ertürk (1998) eğitimi, “bireyin davranışlarında kendi yaşantısı yoluyla ve kasıtlı olarak istendik değişiklik meydana getirme süreci” olarak tanımlamaktadır (Ertürk, 1998, 12).Bilimi çeşitli şekillerde tanımlamak mümkündür. Örneğin, Yıldırım (1991) bilimi “dünyamızda olup biten olguları ‘betimleme ve açıklama’ yoluyla anlama girişimidir” şeklinde tanımlıyor (s. 95).

Cümle arasındaki kısa alıntı örneği;

Bu sonuçları yorumlarken, Robbins ve ark. (2003) “terapistler, bırakılma durumda istemeden de olsa, ergenin ihtiyaçlarını ya da kaygılarını tam olarak karşılamayabilirler” (s. 541) diyerek katkıda bulundular.

Cümle sonundaki kısa alıntı örneği;

Bu konuda kafa karışıklığına sebep olan şey, palyatif bakım sırasındaki rollerin doğasıdır, “tıbbi gereksinimler tıp alanındaki görevlilerce karşılanır; tıbbi olmayan ihtiyaçlar ise ekipteki herhangi birisi tarafından karşılanabilir” (Csikai ve Chaitin, 2006, s. 112).

\section{Dipnotlar}
Tezlerde içeriği genişletici, güçlendirici veya ilave nitelikteki bilgiler (içerik dipnotu) kullanılabilir\footnote{Dipnotlar ile kaynak gösterimi yapılmaz. Dipnotlar tez içerisinde içeriği genişletici, güçlendirci veya ilave nitelikteki bilgileri vermek için kullanılır. Verilen genişletici, güçlendirci veya ilave nitelikteki bilgiler zorunlulukla kaynak içeriyorsa bu kaynak mutlaka kaynaklar bölümünde verilmelidir.}.

Tezlerde içeriği genişletici, güçlendirici veya ilave nitelikteki bilgiler (içerik dipnotu) kullanılabilir.

Dipnot numaraları alıntının hemen sonuna koyulur. Alıntı paragrafsa dipnot numarası paragrafın son kelimesinin üzerine, alıntı bir kavram veya isimse, kavram veya ismin hemen üzerine yazılır.

Metin içerisindeki dipnot numarası; satır hizasının üzerinde1 şeklinde görünür olmalıdır. Numara sonrasında herhangi bir noktalama işareti konmamalıdır.

Dipnot, ilgili sayfanın altına metinden 2 karakter küçük (10 punto) yazı ile yazılmalıdır.

Dipnot çizgisi ile dipnot numarası arasında bir aralık; dipnot numarası ile dipnotun ilk satırı arasında ise yarım aralık bırakılmalıdır. Dipnotlar metinden ince yatay bir çizgi ile ayrılmalıdır.

\section{İkinci Derece Başlık Nasıl: İlk Harfler Büyük}
Lorem ipsum dolor sit amet, consetetur sadipscing elitr, sed diam nonumy eirmod tempor invidunt ut labore et dolore magna aliquyam erat, sed diam voluptua. At vero eos et accusam et justo duo dolores et ea rebum. Stet clita kasd gub rgren, no sea

\subsection{Üçüncü derece başlık nasıl: ilk harf büyük diğerleri küçük}
Lorem ipsum dolor sit amet, consetetur sadipscing elitr, sed diam nonumy eirmod tempor invidunt ut labore et dolore magna aliquyam erat, sed diam voluptua. At vero eos et accusam et justo duo dolores et ea rebum. Stet clita kasd gub rgren, no sea

\subsubsection{Dördüncü derece başlık nasıl: ilk harf büyük diğerleri küçük}
Stet clita kasd gub rgren, no sea takimata sanctus est Lorem ipsum dolor sit amet, consetetur sadipscing elitr, sed diam nonumy eirmod tempor invidunt ut lab ore sit et dolore magna.

\baslikbes{Beşinci derece başlık: dördüncü dereceden sonrası numaralandırılmaz}
% [Normal] sp=after:240
Stet clita kasd gub rgren, no sea takimata sanctus est Lorem ipsum dolor sit amet, consetetur sadipscing elitr, sed diam nonumy eirmod tempor invidunt ut lab ore sit et dolore magna.

\begin{figure}[H]
  \centering
  \includegraphics[width=9.33cm]{sekiller/image6.png}
  \caption{Örnek şekil.}
\end{figure}
Stet clita kasd gub rgren, no sea takimata sanctus est Lorem ipsum dolor sit amet, consetetur sadipscing elitr, sed diam nonumy eirmod tempor invidunt ut lab ore sit et dolore magna.

\begin{table}[H]
  \caption{Tablo örneği.}
  \centering
\begin{tabular}{C{3.05cm}C{2.10cm}C{2.16cm}C{2.52cm}}
\hline\hline
\Ns Kolon A & Kolon B & Kolon C & Kolon D \\
\hline
\Ns Satır A & Satır A & Satır A & Satır A \\
\Ns Satır B & Satır B & Satır B & Satır B \\
\Ns Satır C & Satır C & Satır C & Satır C \\
\hline
\end{tabular}
\end{table}
% [Tablodansonragelenmetin] 
Stet clita kasd gub rgren, no sea takimata sanctus est Lorem ipsum dolor sit amet, consetetur sadipscing elitr, sed diam nonumy eirmod tempor invidunt ut lab ore sit et dolore magna. Stet clita kasd gub rgren, no sea takimata sanctus est Lorem ipsum dolor sit amet, consetetur sadipscing elitr, sed diam nonumy eirmod tempor invidunt ut lab ore sit et dolore magna.

Stet clita kasd gub rgren, no sea takimata sanctus est Lorem ipsum dolor sit amet, consetetur sadipscing elitr, sed diam nonumy eirmod tempor invidunt ut lab ore sit et dolore magna.

